\documentclass[conference]{IEEEtran}

\usepackage{booktabs}
\usepackage{float}
\usepackage{graphicx}
\usepackage{url}

% Working example and scaffold. Replace the italic prompts with your analysis.
% See report/WRITING_GUIDE.md for section purposes and writing examples.
% Keep generated inputs so "all" rebuilds every reported table and figure.

\begin{document}

\title{Warehouse Improvement Decision: Evidence and Recommendation}

\author{\IEEEauthorblockN{Amira de Vries, Luca Jansen, and Noor Bakker}
\IEEEauthorblockA{Warehousing course, University of Twente\\
Group 00 (example names; replace with your group members and number)}}

\maketitle

\begin{abstract}
This example checks an approved S.~P.~Richards inventory file and rebuilds a
zone table and chart in the report. Both describe a stock snapshot; they do
not establish an improvement. For the final report, replace this abstract
with your decision, its measured effect against the current operation, the
main feasibility condition, and the principal uncertainty.
\end{abstract}

\begin{IEEEkeywords}
warehousing, activity profiling, storage assignment, reproducible analysis
\end{IEEEkeywords}

\section{How to use this scaffold}
The final submission is one management report of no more than 12 pages, with
supporting code. Each of the five assessed parts is worth 10 points. You may
combine and reorder the suggested subsections to make one clear argument.
Keep enough evidence in the report and code for every applicable criterion.
Remove this guidance section and the italic prompts before submission.

\subsection{WH-AS1: diagnosis and activity profiling}
The rubric weights the warehouse heatmap at 15\%; two customer-order, two
item-activity, and two inventory profiles at 10\% each; and one time, two
quality, and two productivity KPIs at 5\% each. The formative WH-AS1
submission must address every profile and KPI. The final report can select
the results most relevant to its recommendation, while the code retains
the complete evidence. The course rules also require each student's two
assigned peer reviews for that student's WH-AS1 score.

\subsection{WH-AS2--WH-AS5: five criteria per part}
Each criterion is worth two points. Complete, correct, decision-relevant,
reproducible evidence earns two; a material omission or minor error earns
one; missing or materially incorrect evidence earns zero. The sections below
name the evidence expected for storage, layout, picking, and their combined
trade-offs. The current course rubric controls the final assessment.

\section{Decision and case}
\subsection{Proposed decision and expected effect}
The S.~P.~Richards Philadelphia DC case supplies the operational data for this
project~\cite{spr}. This example establishes a reproducible inventory baseline.
It makes no storage or picking recommendation.

\emph{Student prompt: State the management decision first. Identify the area,
the proposed change, the current practice, and the expected operational
effect. Give the conditions under which the recommendation remains feasible.}

\subsection{Operating context}
\emph{Student prompt: Define the warehouse area, relevant product and order
flows, and the decision boundary. Describe the present operation briefly so
the comparison has a clear baseline.}

\section{Diagnosis and current operation}
\subsection{Data scope and checks}
The pipeline reads the \texttt{DC23ACTIVE} sheet from the approved archive. It
counts active-location records and distinct SKUs by zone, then sums recorded
units on hand. The source archive and extracted file are checked against the
SHA-256 values in \texttt{data/sources.json}; the originals are left unchanged.

\emph{Student prompt: For each decisive measure, give the source file, unit,
time period, calculation, and any missing or assumed information. Make the
checks repeatable in the code.}

\subsection{Activity profiles and KPIs}
\emph{Student prompt: Explain the heatmap, order, item, and inventory patterns
that motivate the decision. Report the time, quality, and productivity KPIs
with their denominators. Keep the full WH-AS1 evidence in the code even if
the final report shows only selected results.}

\subsection{A generated baseline table and chart}
\input{artifacts/summary.tex}
\input{artifacts/zone_chart_count.tex}

\begin{table}[H]
\caption{Active-location inventory by zone in the DC23ACTIVE snapshot
($n=\activeLocationCount$ location records; units on hand are selling units).}
\label{tab:zone-summary}
\centering
\footnotesize
\setlength{\tabcolsep}{3pt}
\input{artifacts/zone_table.tex}
\end{table}

\begin{figure}[H]
\centering
\includegraphics[width=\columnwidth]{artifacts/zone_locations.png}
\caption{Active-location records by zone in the DC23ACTIVE snapshot
($n=\activeLocationCount$ records; bar values count location records).}
\label{fig:zone-locations}
\end{figure}

Table~\ref{tab:zone-summary} and Fig.~\ref{fig:zone-locations} describe one
stock snapshot. The chart sorts zones by active-location count, while the
table also reports SKUs and units. Neither measures demand, picking time, or
storage capacity. Those quantities need separate evidence before they can
support a policy choice.

\section{Storage assignment and space (WH-AS2)}
\subsection{Scope, data, and assumptions}
\emph{Student prompt: Define products, locations, flows, resources,
constraints, and the performance measure. Trace each model parameter to
supplied data or an explicit assumption. Document evidence gaps.}

\subsection{Assignment logic and storage policy}
\emph{Student prompt: State the assignment model compactly, including its
objective and key constraints. Explain how product and demand characteristics
lead to a policy that staff can implement. Keep full matrices in the code.}

\subsection{Baseline comparison and feasibility}
\emph{Student prompt: Compare the proposal with current assignments using
absolute and relative effects. Check capacity, compatibility, replenishment,
and service. Explain the main sensitivity and what management must monitor.}

\section{Layout design (WH-AS3)}
\subsection{Requirements, departments, and flows}
\emph{Student prompt: Quantify space, storage, throughput, adjacency, and
handling needs. Show the proposed departments, interfaces, and flows in a
labelled layout figure. Connect the arrangement to the storage policy.}

\subsection{Quantitative and operational checks}
\emph{Student prompt: Calculate distances, space use, congestion proxies, or
other suitable measures against the current layout. Check access, safety,
replenishment, staging, and equipment constraints. Explain why this layout is
preferable to a credible alternative and what implementation would require.}

\section{Picking design (WH-AS4)}
\subsection{Process and design logic}
\emph{Student prompt: Define the picking process, staff roles, release rule,
and operating conditions. Justify routing, batching, zoning, and storage
interactions using the observed workload.}

\subsection{Matched performance and feasibility}
\emph{Student prompt: Test the picking change against current practice under
matched demand. Report the effect on decision-relevant measures. Check labour,
equipment, replenishment, service, variability, and exception handling.
Explain the expected benefit, burden, risks, and boundary conditions.}

\section{Integrated trade-offs (WH-AS5)}
\subsection{Fair comparison and performance}
\emph{Student prompt: Compare storage, layout, and picking options using the
same baseline, demand, units, and constraints. Quantify effects on cost,
throughput, space, service, workload, and risk where evidence permits.}

\subsection{Interactions, uncertainty, and limits}
\emph{Student prompt: Show where interventions reinforce each other, conflict,
or shift work elsewhere. Vary a material assumption and show whether the
preferred option changes. Bound any result that the data cannot support.}

\subsection{Recommendation, sequence, and monitoring}
\emph{Student prompt: Give management a prioritized action, rollout sequence,
conditions for implementation, and indicators to monitor. State what result
would cause you to revise the decision.}

\section{Reproducing the results}
From the project directory, run \texttt{run.cmd all} on Windows or
\texttt{./run.sh all} on macOS or Linux. Docker installs Python and LaTeX. The
command checks original files under \texttt{data/raw}, writes generated tables
and the chart under \texttt{artifacts}, and builds
\texttt{artifacts/main.pdf}. Keeping source data, code, and generated outputs
separate follows reproducible research guidance~\cite{turingway}.

\bibliographystyle{IEEEtran}
\bibliography{report/references}

\end{document}
